\section{Crop Envelopes, Rotation Constraints, and Optimization}
\label{sec:optimization}

\subsection{Agronomic Crop Catalog and Biophysical Filtering}
\texttt{MASTER\_CROPS\_CATALOG} defines twelve representative field, forage, horticultural, and bio-drilling cover crops spanning seven botanical families (Poaceae, Fabaceae, Asteraceae, Brassicaceae, Solanaceae, Amaranthaceae). Before exhaustive sequence evaluation, \texttt{PlanOptimizerEngine} invokes \texttt{\_filter\_candidate\_crops} to prune inadmissible species using field constraints and climate observations:
\begin{enumerate}[leftmargin=1.5em]
  \item \textbf{Farm Covenants \& Infrastructure}: Blacklisted species are excluded; when at least three whitelisted crops are provided, search narrows to the whitelist. Crops requiring protected cultivation are filtered unless greenhouse infrastructure is flagged, and low-spot fields with elevated waterlogging exposure retain only crops with sufficient flood tolerance ($\ge 2.0$).
  \item \textbf{Thermal Envelope \& Frost/Heat Pruning}: Each crop carries an optimal thermal window $[T_{\mathrm{opt,min}}, T_{\mathrm{opt,max}}]$ and lethal frost threshold $T_{\mathrm{frost,kill}}$. Warm-season and cool-season crops are tested against observed seasonal temperatures, frost-day counts (\texttt{frost\_days\_count}), coldest night temperature, and heatwave counts (\texttt{heatwave\_days\_count}) from NASA POWER.
  \item \textbf{Edaphic $\mathrm{pH}$ Lockout Pruning}: Soil $\mathrm{pH}$ resolved from local tests or SoilGrids must lie within $\pm 1.6~\mathrm{pH}$ units of $[\mathrm{pH}_{\mathrm{opt,min}}, \mathrm{pH}_{\mathrm{opt,max}}]$; otherwise the crop is excluded from the feasible candidate set (falling back gracefully only if filtering leaves fewer than three candidates).
\end{enumerate}

\subsection{Rotational Break Rules and Sequential Carryover}
Every ordered 3-year triplet $(c_1, c_2, c_3)$ is simulated sequentially across Years $1 \to 2 \to 3$, carrying forward updated state $(C_t, N_{\mathrm{min},t}, \mathrm{PSI}_t, \mathrm{family}_{t-1})$ and preceding legume/scavenger credits. Agronomic sequence constraints enforce:
\begin{itemize}[leftmargin=1.5em]
  \item \textbf{Phytosanitary Break Intervals} (\texttt{rotational\_break\_needed\_years}): Solanaceae crops (\texttt{processing\_tomatoes}, \texttt{bell\_peppers}) requiring a $\ge 3$-year soil-borne disease break are limited to \emph{at most once} in any 3-year sequence and check any reported Year~$0$ preceding crop. Crops with $\ge 2$-year breaks prohibit consecutive same-family succession, and consecutive pure cover-crop seasons (\texttt{daikon\_radish\_cover}) are disallowed.
  \item \textbf{Continuous-Family Pathology Drag}: When adjacent crops share a botanical family without meeting ideal break length, expected yield suffers a $12\%$ continuous-family drag ($0.88$ multiplier), avoiding arbitrary fabricated dollar fines while reflecting agronomic yield depression.
\end{itemize}

\subsection{Composite Yield Attenuation and Economic Valuation}
For each year $t \in \{1, 2, 3\}$, baseline potential yield $Y_{0,t}$ (calibrated against FAOSTAT national production yields when available for the parcel's country) is scaled by multiplicative abiotic and biological stress factors:
\begin{equation}
  Y_t = Y_{0,t} \cdot \mu_{\mathrm{water}} \cdot \mu_{\mathrm{pH}} \cdot \mu_{\mathrm{thermal}} \cdot \mu_{\mathrm{compaction}} \cdot \mu_{\mathrm{family}},
\end{equation}
where:
\begin{itemize}[leftmargin=1.5em]
  \item $\mu_{\mathrm{water}} = \max(0.35,\, 1 - \Delta Y_{\mathrm{water}})$ derives directly from FAO-33 stress $K_y(1 - K_s)$ (Section~\ref{sec:hydrology}).
  \item $\mu_{\mathrm{pH}}$ applies a nutrient lockout attenuation of $0.14$ per $\mathrm{pH}$ unit outside $[\mathrm{pH}_{\mathrm{opt,min}}, \mathrm{pH}_{\mathrm{opt,max}}]$.
  \item $\mu_{\mathrm{thermal}}$ accounts for seasonal departure from $[T_{\mathrm{opt,min}}, T_{\mathrm{opt,max}}]$, acute NASA POWER heatwave exposure, and live ESP32 canopy ambient temperature departures when sensor telemetry is present.
  \item $\mu_{\mathrm{compaction}}$ attenuates root-limited yield when $\mathrm{PSI}_t$ exceeds the crop's critical penetrometer threshold.
\end{itemize}

Farm-gate revenue prices $Y_t$ using market benchmarks ingested from World Bank, FAOSTAT, and Eurostat indices~\cite{faostat2024} (falling back to catalog defaults only when external series are unavailable for that commodity). Synthetic fertilizer costs scale linearly with simulated nitrogen demand $N_{\mathrm{synth}}$, phosphorus removal, and potassium removal priced against World Bank urea, TSP, and KCl nutrient tariffs.

\subsection{Five-Strategy Combinatorial Selection}
\texttt{PlanOptimizerEngine.optimize\_all\_strategies} evaluates every admissible 3-year triplet and selects five non-redundant strategic champions suited to distinct farmer priorities:
\begin{enumerate}[leftmargin=1.5em]
  \item \textbf{Most Income (\texttt{most\_income})}: Maximizes cumulative 3-year net farm margin ($\mathrm{USD}$).
  \item \textbf{Best Soil Health (\texttt{best\_soil\_health})}: Maximizes weighted multi-year $\Delta\mathrm{SOC}$, compaction reduction ($\Delta\mathrm{PSI} < 0$), erosion suppression ($\mathrm{RUSLE}~E$), and biological nitrogen enrichment.
  \item \textbf{Least Effort (\texttt{least\_effort})}: Minimizes field labor hours, machinery passes, and irrigation pumping burden while preserving positive farm margin.
  \item \textbf{Highest ROI (\texttt{highest\_roi})}: Maximizes capital return ratio ($\text{Net Profit} / \text{Operating Cost}$), favoring low-input drought-hardy or legume-subsidized sequences.
  \item \textbf{Least Crop Loss Risk (\texttt{least\_crop\_loss\_risk})}: Minimizes composite multi-hazard risk index spanning unmet water deficit, observed frost/heat extreme vulnerability, continuous-family disease exposure, and soil compaction excess.
\end{enumerate}
